\scanpage{52}
Je travaillerai \deleted{donc} avec l'iso.
\[
(23)\qquad P\Sigma\simeq V^*
\]
plutôt qu'avec $P\Sigma\simeq V(\underline\omega)^*$, \deleted{\unsure{ce qui}} qui respecte l'action de $\mathbb C^*$ \deleted{\unsure{\illegible{1}}} (au lieu de la remplacer par l'action inverse), mais qui, par contre, n'est pas compatible aux projections --- on a seulement \unsure{d'int.}
\diagram{\begin{tikzcd}[ampersand replacement=\&, row sep=large,column sep=large]
P\Sigma_{\unsure{\Sigma}}\arrow[r,"\sim","(23)"']\arrow[d,"{\mathrm{proj.can.}}"']\&V^*\arrow[d,"{\mathrm{proj.can.}}"]\\
\Sigma\arrow[r,"\sigma"]\&\Sigma\simeq V^*/\mathbb R^*
\end{tikzcd}}
\margin{(24)}

\margin{Toute la digression précédente est un pur ``\unsure{sexe des anges}''. À retenir : pour le fibré vectoriel \deleted{\unsure{F}} \unsure{intrinsèque} $\underline L$ \deleted{\unsure{pour}} sur $\Sigma$ (\unsure{au moins dans le cas $n=3$}) au dessus d'un point $R_i\in\Sigma$, il y a un point $\widetilde R_i$ canonique (\unsure{impliquant des choix} \deleted{\unsure{amples}} \unsure{identifiés à} un $\omega\in\underline\omega$), et qui \deleted{$\unsure{\Gamma(\Sigma,L)}$} \deleted{$\unsure{\Gamma(\Sigma,L)}$} \unsure{identifie les} $\check R_i$ à des éléments de $\Gamma(\Sigma,\underline L)(\omega)\overset{\mathrm{déf}}=\check V$, dont la somme est $0$, et qui engendrent, \unsure{d'où un isom} $V\simeq\mathcal O^J/\mathcal O$ (diagonale) \unsure{--- le tout} \unsure{avec les flèches} \deleted{\unsure{et $\omega\in$}} \uline{\unsure{sur $\Sigma$}}. Ceci \unsure{se relativise} sur $\mathbb Z$. $\Sigma$ : un \deleted{\unsure{$V^*/\mathbb C^*$}} $V^*/\mathbb C^*$, \unsure{désigne par} $\check V$, $P\Sigma$ \unsure{``précisément comme les droites de la droite $L_x$''}, \uline{avant} de l'opération des $\mathbb C^*$) et \unsure{n'y a pas} dans $V=\Gamma(\Sigma,\underline L)(\omega)$ \unsure{qu'il faut prendre} (\unsure{plutôt} $V=\operatorname{Ker}(\mathcal O^J\to\mathcal O)$), $V=\mathcal O^J/\mathcal O$. Cet \uline{iso} est bien précis sur les \unsure{3 isom}, on a \unsure{une} commut\unsure{$^{\mathrm{n}}$} à l'action de $D\,(\simeq\unsure{\mathfrak S_3})$.}

Notons en passant que si $\tau$ désigne la conjugaison complexe (de $V_{\mathbb C}$ ou $\Sigma$) on aura
\[
(25)\qquad\tau\deleted{\unsure{\circ}}\sigma=\underline a\quad\text{ou ``antipodisme'' de }\Sigma\quad\ldots
\]
Je m'intéresse à l'ouvert
\[
(26)\qquad\Sigma^*=\Sigma-\{R_i\mid i\in J\}
\]
et à : \deleted{\unsure{B}}
\[
(27)\qquad P\Sigma^*=P\Sigma\mid\Sigma^*\overset{(23)}\simeq V^{**}=V-\bigcup_{i\in J}\Delta_i
\]
où les $\Delta_i$ sont les $n=m/2$ \uline{diagonales} du
\mspage{51}
